\section{Adaptive spectral rounding and the Hamiltonicity consumer}
\label{sec:rounding}
We first recall the elementary decomposition used at each step.
If \(T\) is a doubly stochastic matrix, its positive support has
a perfect matching: a set of rows \(U\) sends mass \(|U|\)
into columns of total capacity \(|N(U)|\), so Hall's condition holds.
Choose such a permutation matrix \(M\) and let \(w>0\) be the
smallest matched entry.
Then
\[
 T=wM+(1-w)\frac{T-wM}{1-w}.
\]
If \(w=1\) the decomposition is complete; otherwise the second
matrix is doubly stochastic and has fewer positive entries.
Iteration gives a finite convex combination of permutation
matrices. In symmetric bipartite block form, these are adjacency
matrices of perfect matchings. In particular, for a stochastic
bipartite kernel \(P_i\) there is a random perfect matching \(M_i\)
with \(\E M_i=P_i\). This decomposition is an established input,
not a new claim.

\paragraph{Formal finite matching selector.}
The classical finite selection input has now been proved in Lean in
a form that can also drive a deterministic moment argument.
For every nonnegative doubly stochastic cross matrix \(P\) and
every real cost matrix \(C\), there is a permutation \(f\) with
\(P_{i,f(i)}>0\) for every \(i\) such that its permutation matrix
\(M_f\) satisfies the signed inequality
\[
 \sum_{i,j} C_{ij}\bigl((M_f)_{ij}-P_{ij}\bigr)\leq0.
\]
For completeness, choose a supported matching using the exact Hall
mass count above.  Let \(w>0\) be its smallest matched entry.
If \(w=1\), stochasticity forces \(P=M_f\).
Otherwise
\[
 R=\frac{P-wM_f}{1-w},\qquad P=wM_f+(1-w)R
\]
is nonnegative and doubly stochastic, its positive support is contained
in that of \(P\), and at least one matched entry has disappeared.
Induct on the number of positive entries.
If \(M_f\) already has cost at most that of \(P\), use it.
Otherwise the displayed convex identity forces \(R\) to have
strictly smaller cost than \(P\), and the induction hypothesis on
\(R\) supplies a supported matching with no larger cost.
This proves the signed conclusion without a probability assumption.
The maintained owners are \texttt{HamiltonianStochasticMatching},
\texttt{HamiltonianStochasticPeeling} and
\texttt{HamiltonianMatchingCostSelection}.
This is the established finite assignment principle; no originality
is claimed for it.  The full one-step matrix-moment estimate, finite adaptive
construction and same-graph normalized spectral and observable
readouts are now proved in Lean, as detailed below.

\paragraph{Formal trace-power estimate for deterministic selection.}
The next analytic input is now also proved in Lean, uniformly in the
matrix order and the moment order.
For real symmetric \(A,X\), put \(A_t=A+tX\) and
\(g_p(t)=\operatorname{tr}(A_t^p)\), where \(p\geq2\) is an integer.
Scalar differentiation of every matrix entry gives the exact formulas
\[
 g_p'(t)=p\operatorname{tr}(A_t^{p-1}X),\qquad
 g_p''(t)=p\sum_{j=0}^{p-2}
          \operatorname{tr}(A_t^{p-2-j}XA_t^jX).
\]
The matrix power direction retains every noncommuting word.
Diagonalize \(A_t\) only after differentiating:
write \(A_t=U\diag(\lambda_r)U^\top\) and \(B=U^\top XU\).
No derivative of the eigenbasis is used.
Every mixed trace is exactly the full ordered-pair sum
\[
 \operatorname{tr}(A_t^aXA_t^bX)
   =\sum_{r,s}B_{rs}^{\,2}\lambda_r^a\lambda_s^b.
\]
In particular all signed eigenvalues and return pairs remain present.
The actual operator norm gives both budgets
\[
 \sum_s B_{rs}^{\,2}\leq\|X\|^2,\qquad
 \sum_r B_{rs}^{\,2}\leq\|X\|^2.
\]
For \(x,y\geq0\) and \(0\leq j\leq n\),
\(x^jy^{n-j}\leq x^n+y^n\).
Apply this with \(n=p-2\) to each complete mixed trace,
then sum using the two square-mass budgets.  This proves
\[
 \boxed{\quad
 g_p''(t)\leq
 2p(p-1)\|X\|^2\sum_r|\lambda_r(A_t)|^{p-2}.
 \quad}
\]
The bound uses the eigenvalue mass instead of replacing it by the
ambient order times a maximum entry.
The maintained proof consists of
\texttt{HamiltonianTracePowerScalars},
\texttt{HamiltonianTraceSpectralBasis},
\texttt{HamiltonianTracePowerDerivative} and
\texttt{HamiltonianTraceHessianBudget}.
It is a supporting classical matrix estimate; no originality is
claimed for the estimate in isolation.
The eigenvalue-mass interpolation and the complete regularized
squared-moment step are now also proved in Lean.
For every finite nonnegative family \(a_r\) and \(0\leq q\leq1\),
\[
 \sum_r a_r^q\leq n^{1-q}\left(\sum_r a_r\right)^q.
\]
The proof raises the power-mean inequality with exponent \(1/q\)
back to the power \(q\), treating \(q=0\) separately.
It includes empty finite sets.
Apply it to \(a_r=|\lambda_r|^p\) and \(q=1-2/p\).
For even \(p\geq2\), the exact Hermitian trace identity gives
\[
 \sum_r|\lambda_r|^{p-2}
 \leq n^{2/p}\bigl(\operatorname{tr}A^p\bigr)^{1-2/p},
 \qquad \operatorname{tr}A^p\geq0.
\]
Now write \(g(t)=\operatorname{tr}(A+tX)^p\),
\(q=2/p\), and \(h(t)=g(t)+\varepsilon>0\), where
\(\varepsilon>0\).
The full differentiated expression is
\[
 \frac{d^2}{dt^2}h(t)^q
 =g''(t)q h(t)^{q-1}
  +(g'(t))^2q(q-1)h(t)^{q-2}.
\]
The second term is nonpositive, since \(0<q\leq1\).
In the first term the Hessian bound and the interpolation above give
\(g''(t)\leq K g(t)^{1-q}\), with
\(K=2p(p-1)n^{2/p}\|X\|^2\).
Using \(0\leq g(t)\leq h(t)\), the exact power cancellation
\(h^{1-q}h^{q-1}=1\) yields
\[
 \frac{d^2}{dt^2}h(t)^q
 \leq4(p-1)n^{2/p}\|X\|^2.
\]
This is uniform in \(A\) and \(\varepsilon>0\).
The one-sided Taylor theorem is proved by applying derivative
monotonicity first to \(f'(t)-Bt\), then to
\(f(t)-f'(0)t-Bt^2/2\); it needs an upper bound on \(f''\),
rather than an absolute second-derivative bound.
Consequently, if
\(\operatorname{tr}(A^{p-1}X)\leq0\), then
\[
 \boxed{\quad
 \bigl(\operatorname{tr}((A+X)^p)+\varepsilon\bigr)^{2/p}
 \leq
 \bigl(\operatorname{tr}A^p+\varepsilon\bigr)^{2/p}
 +2(p-1)n^{2/p}\|X\|^2 .
 \quad}
\]
These proofs are maintained in
\texttt{HamiltonianPowerMassInterpolation},
\texttt{HamiltonianTraceMomentMass},
\texttt{HamiltonianRegularizedPower},
\texttt{HamiltonianSquaredTraceCurvature},
\texttt{HamiltonianOneSidedTaylor} and
\texttt{HamiltonianSquaredTraceStep}.
They are classical analytic supporting estimates.
The complete analytic finite-path iteration and its operator-norm
readout are now also proved.
Suppose \(Z_0=0\), \(Z_{i+1}=Z_i+X_i\), the matrices are
Hermitian, \(\|X_i\|\leq2\), and
\(\operatorname{tr}(Z_i^{p-1}X_i)\leq0\) for every \(i<d\).
The previous step, summed with its full constants, gives
\[
 \bigl(\operatorname{tr}(Z_d^p)+\varepsilon\bigr)^{2/p}
 \leq \varepsilon^{2/p}+8(p-1)d\,n^{2/p}.
\]
To read the actual operator norm, the orthogonal spectral
decomposition gives
\(\|A\|\leq\|(\lambda_r)_r\|_\infty\), with the norm of the
empty vector equal to zero.
Every nonnegative even eigenvalue power is bounded by the full
trace sum, so
\[
 \|A\|^2\leq\bigl(\operatorname{tr}(A^p)\bigr)^{2/p}.
\]
Both conclusions include the empty finite index type.
Taking \(\varepsilon=1\), rather than discarding the initial
regularization, yields the compiled finite-path estimate
\[
 \boxed{\quad \|Z_d\|^2\leq1+8(p-1)d\,n^{2/p}.\quad}
\]
The two additional maintained owners are
\texttt{HamiltonianTraceNormReadout} and
\texttt{HamiltonianFiniteTracePath}.
This is the actual operator norm, with the dimension occurring only
as \(n^{2/p}\).
The path theorem explicitly requires the signed direction condition
at every step; no conditional concentration hypothesis is used.
The actual local matching step for the joint packet is now also
proved in Lean.
For a cross matrix \(C\), let
\[
 B(C)=\begin{pmatrix}0&C\\C^\top&0\end{pmatrix},\qquad
 {\cal L}(C)=\operatorname{diag}
 \bigl(B(C),\diag(B(C)\one_{S_1}),\diag(B(C)\one_{S_2})\bigr).
\]
Here \(S_1,S_2\) are arbitrary subsets of the full bipartite vertex set.
The packet has order \(n=3N=6m\) for an \(m\)-by-\(m\) cross matrix.
Its first block is exactly \(B(C)\), and its other two block
diagonals are exactly the requested row observables.
The map \({\cal L}\) is real linear and has Hermitian values.
If \(P\) is nonnegative and doubly stochastic, every row of \(B(P)\)
has absolute mass one.
Each observable diagonal entry is bounded in modulus by that same
row mass.
The established symmetric row-budget theorem therefore gives
\[
 \|{\cal L}(P)\|\leq1.
\]
The same proof controls the lift by any uniform absolute row-mass
budget; it uses the full repaired-kernel subset estimate.

For any current Hermitian packet \(Z\), the complete trace
functional
\(\phi(C)=\operatorname{tr}(Z^{p-1}{\cal L}(C))\)
is linear.
With \(E_{ij}\) the single-entry matrix units, its exact finite
coordinate representation is
\[
 \phi(C)=\sum_{i,j}\phi(E_{ij})C_{ij}.
\]
Thus the already proved arbitrary signed-cost matching selector
produces a supported permutation matrix \(M\) with
\[
 X={\cal L}(M-P),\qquad
 X=X^\top,\qquad \|X\|\leq2,\qquad
 \operatorname{tr}(Z^{p-1}X)\leq0.
\]
The norm bound uses linearity and the two actual stochastic-kernel
bounds \(\|{\cal L}(M)\|,\|{\cal L}(P)\|\leq1\).
Every hypothesis needed by the finite moment path theorem has
therefore been discharged for this actual local supported matching
step, for the same unrestricted \(S_1,S_2\).
The maintained owners are
\texttt{HamiltonianLinearMatchingDirection},
\texttt{HamiltonianJointObservableLift} and
\texttt{HamiltonianBipartiteRoundingStep}.
These are classical supporting constructions.
The exact repair-and-freshness step is now also proved in Lean.
Suppose the current cross adjacency matrix \(H\) has entries zero
or one, every row and column sums to \(i\), and it is supported
on the original \(P\).
Define the actual deleted mass entrywise by
\[
 Q_{ab}=P_{ab}H_{ab}.
\]
If \(0\leq P_{ab}\leq\alpha\), then
\[
 0\leq Q\leq P,\qquad
 \sum_b Q_{ab}\leq i\alpha,\qquad
 \sum_a Q_{ab}\leq i\alpha.
\]
These are exact mask bounds, with no independence or disjoint
pairing assumption.
Under \(i\alpha L^2\leq c_\delta\), apply the full repair theorem
to the original \(P\), yielding a stochastic \(W\) with
\[
 W_{ab}=0\ \Longleftrightarrow\
 (P-Q)_{ab}=P_{ab}(1-H_{ab})=0,\qquad
 \|{\cal L}(W-P)\|\leq C_\delta^{\rm obs}i\alpha L.
\]
Select the actual joint signed matching using \(W\).
Its positive support forces \(H_{a,f(a)}=0\) and
\(P_{a,f(a)}>0\).
Thus \(H+M_f\) is still binary, is supported on \(P\), and
has every row and column sum exactly \(i+1\).
The same witness supplies all stochastic, support, drift and
joint-moment direction properties.
Every repair uses the original kernel and its original gap;
the argument does not repeatedly halve the gap of successive repairs.
The maintained owners are \texttt{HamiltonianUsedEdgeMass} and
\texttt{HamiltonianUsedEdgeJointStep}.
The complete finite-prefix induction and its quantitative error are
now also proved in Lean.
The maintained prefix records the actual binary supported \(H_i\),
the whole centered matrix path, every Hermitian increment and signed
direction, and a drift matrix satisfying
\[
 {\cal L}(H_i-iP)=Z_i+R_i,\qquad
 \|R_i\|\leq C_\delta^{\rm obs}\alpha L\,i^2 .
\]
At extension, append \(X_i={\cal L}(M_{i+1}-W_i)\) and set
\[
 Z_{i+1}=Z_i+X_i,\qquad
 R_{i+1}=R_i+{\cal L}(W_i-P).
\]
Every earlier path entry is retained.
The exact linear matrix identity proves the new error identity.
The full repair bound gives
\[
 \|R_{i+1}\|
 \leq C_\delta^{\rm obs}\alpha L(i^2+i)
 \leq C_\delta^{\rm obs}\alpha L(i+1)^2.
\]
The global smallness hypothesis implies each prefix repair's
smallness hypothesis.
Fix an even \(p\geq2\) before this construction.
Induction starts with zero adjacency, centered path and drift.
It constructs an actual final binary supported cross matrix
\(H=H_p\), with every row and column sum \(d\).
The witness may depend on the chosen moment order.
Applying the compiled path moment estimate to this same \(H_p\)
gives
\[
 \boxed{\quad
 \|{\cal L}(H-dP)\|
 \leq\sqrt{1+8(p-1)d(6m)^{2/p}}
       +C_\delta^{\rm obs}\alpha Ld^2.
 \quad}
\]
The owners are \texttt{HamiltonianRoundingPrefix},
\texttt{HamiltonianRoundingPrefixExtension} and
\texttt{HamiltonianRoundingConstruction}.
The full actual finite construction is proved, rather than assumed
as a favorable matrix path.
The logarithmic moment choice and the final same-graph readouts are
now also proved in Lean.
Put \(n=6m\) and choose the even integer
\[
 p=2\lceil2\log n\rceil.
\]
For \(m\geq1\), the scalar proof gives
\[
 p\geq2,\qquad p-1\leq5\log n,\qquad
 n^{2/p}\leq2,\qquad
 1\leq L=\log(4m),\quad \log n\leq2L.
\]
Consequently, for every integer \(d\geq1\),
\[
 1+8(p-1)d\,n^{2/p}\leq256dL.
\]
The initial unit is included: \(dL\geq1\) absorbs it in this
explicit bound.
For the actual matrix constructed at this chosen \(p\), division
by the proved positive \(d\) yields
\[
 \boxed{\quad
 \|{\cal L}(H/d-P)\|
 \leq16\sqrt{L/d}+C_\delta^{\rm obs}d\alpha L.
 \quad}
\]
The normalization uses the exact identity
\(H/d-P=(H-dP)/d\), with its nonzero denominator checked.
Compression onto the first block controls the full spectral norm.
The other two diagonal blocks control the signed row sums for the
same unrestricted subsets \(S_1,S_2\), on this same matrix.

Finally set \(E=\{(a,b):H_{ab}=1\}\).
The binary-entry proof identifies its genuine adjacency matrix
with \(B(H)\); the row and column sum identities give exactly
\(d\) distinct incident edges at each vertex.
Every edge lies in the positive support of the original \(P\).
For every finite vertex subset \(S\), the adjacency sum is exactly
\(\deg_E(v,S)\), proved by a finite indicator count.
Thus both observables are actual neighbor counts of one simple
regular bipartite graph, not separate matrix witnesses.
The additional maintained owners are
\texttt{HamiltonianLogMomentChoice},
\texttt{HamiltonianLogarithmicRounding},
\texttt{HamiltonianNormalizedRounding},
\texttt{HamiltonianRoundingReadout} and
\texttt{HamiltonianRoundedEdgeSet}.
The last owner also proves positivity of the dimension-independent
constants
\[
 c_\delta=c_\delta^{\rm repair},\qquad
 C_\delta=\max\{16,C_\delta^{\rm obs}\},
\]
and the complete advertised bound (1), simultaneously with both
neighbor-count conclusions.
In fact the displayed stronger estimate does not require the
additional \(L/d\) term in (1).

\begin{proof}[Proof of Theorem~\ref{thm:rounding}]
Apply the complete deterministic joint construction above, with any
two specified subsets (or with both empty for the spectral assertion).
The full signed prefix identity and every repair are taken relative
to the original kernel and its original gap.
At the even moment \(p=2\lceil2\log(6m)\rceil\), the explicit
scalar estimate and positive-degree normalization give
\[
 \left\|\frac{A_H}{d}-P\right\|
 \leq16\sqrt{L/d}+C_\delta^{\rm obs}d\alpha L.
\]
The binary support and exact degree identities give a simple
\(d\)-regular spanning graph.
Take \(c_\delta=c_\delta^{\rm repair}>0\) and
\(C_\delta=\max\{16,C_\delta^{\rm obs}\}>0\).
Both positive terms are bounded by their corresponding terms
with coefficient \(C_\delta\); adding the nonnegative \(C_\delta L/d\)
gives (1).
All constants depend only on \(\delta\).
This proof is the maintained Lean proof; it requires no probabilistic
concentration input.
\end{proof}

\begin{lemma}[Rounding preserves two specified row observables]
\label{lem:observablerounding}
Under the hypotheses of Theorem~\ref{thm:rounding}, fix any
two vertex subsets \(S_1,S_2\).
The graph \(H\) can be chosen to satisfy its operator bound
and, simultaneously,
\[
 \max_{\substack{v\\S\in\{S_1,S_2\}}}
 \left|\frac{\deg_H(v,S)}d-(P\one_S)(v)\right|
 \leq C_\delta\left(d\alpha L+\sqrt{L/d}+L/d\right).
\]
There is no size restriction on these subsets.
\end{lemma}
\begin{proof}
Use the same joint construction with these \(S_1,S_2\).
The corresponding observable blocks have diagonal entries
\[
 \frac{\deg_H(v,S_j)}d-(P\one_{S_j})(v),\qquad j=1,2.
\]
The modulus of each entry is at most the norm of the complete
joint packet.
The explicit normalized bound and the same
\(C_\delta=\max\{16,C_\delta^{\rm obs}\}\) prove the assertion.
The finite edge-set dictionary and both counts use the very same
supported graph as the spectral estimate.
\end{proof}

\paragraph{Formalization status.}
The repaired stochastic witness now has a Lean proof of the complete
row and column absolute-mass estimate, uniformly over every subset
observable.  With \(L=\log(2N)\), its explicit constant is
\[
 C_\delta^{\mathrm{obs}}
 =3C_\delta^{\mathrm{scale}}+\frac1{\log4}.
\]
Indeed, write \(T=P-Q\), \(r=C_\delta^{\mathrm{scale}}\eta L\)
and \(W=\diag(1+z_L)T\diag(1+z_R)\), using the same scaling
witness as in Corollary~\ref{cor:forbiddenrepair}.
The smallness condition gives \(0\leq r\leq1\), and the exact
entrywise difference is
\[
 W_{ij}-P_{ij}
 =(z_L(i)+z_R(j)+z_L(i)z_R(j))T_{ij}-Q_{ij}.
\]
Consequently every row and every column has absolute mass at most
\[
 \eta+2r+r^2\leq\eta+3r
 \leq C_\delta^{\mathrm{obs}}\eta L .
\]
This retains the deleted mass and the quadratic cross term.
Summing over any subset gives the stated predictable observable drift
with no restriction on its size.  The maintained owners are
\texttt{HamiltonianScaledCrossMass} and
\texttt{HamiltonianForbiddenObservableDrift}.
The logarithmic moment choice, full normalization, simple edge-set
construction, spectral readout and both actual neighbor-count readouts
are now all formalized in the five owners listed above.
Theorem~\ref{thm:rounding} and Lemma~\ref{lem:observablerounding}
have complete maintained Lean proofs.
Other headline claims of the paper remain under formalization.


\begin{proof}[Proof of Corollary~\ref{cor:weightedhamilton}]
First suppose \(\alpha\leq N^{-a}\) and take
\(d=\lfloor N^{a/2}\rfloor\).
Then \(d\alpha L^2\to0\), and every term in (1) tends to zero.
Since both \(A_H/d\) and \(P\) have the exact eigenvectors
\(\one,s\) at eigenvalues \(1,-1\), Weyl's inequality on
their common complement gives
\[
 \lambda_2(H)\leq(1-\delta/2)d
\]
for sufficiently large \(N\).
The polynomial degree exceeds
\((6\delta^{-1}\log N)^K\).
Equation (2) makes \(H\), and hence the support of \(P\),
Hamiltonian.

For (3), choose a fixed \(A_\delta>0\) sufficiently large
and set \(d=\lceil A_\delta L^{K+1}\rceil\).
For all sufficiently large \(N\), this is at most \(N/2\)
and exceeds \((6\delta^{-1}\log N)^K\).
Under (3),
\[
 d\alpha L^2\leq 2A_\delta b_\delta,\qquad
 d\alpha L\leq 2A_\delta b_\delta/L.
\]
Choose \(b_\delta\) small enough for the first quantity to
be at most the constant \(c_\delta\) in Theorem~\ref{thm:rounding}.
The second tends to zero, as do the two sampling errors.
The same spectral-gap and Hamiltonicity consumer applies.
The concentration obstruction is its contrapositive.
\end{proof}

\begin{remark}[What is independent and what is quantitative]
The finite rounding bound (1) is proved without a Hamiltonicity
input. It preserves the spectrum, exact degrees, and simplicity
simultaneously. The Hamiltonicity conclusion then uses the
independent recent theorem (2).
The constants \(c_\delta,C_\delta\) are controlled by the
scaling contraction and the deterministic matrix-moment budget; their numerical
values have not been optimized.
The starting order in the Hamiltonicity corollary is
asymptotic and has not been computed.
\end{remark}
